\section{Plate Collision}
\label{sec:release}

When a chunk is released, the gripper opens one of its two plates. The
opening side sweeps the adjacent pallet row: if that row already contains
a box overlapping the chunk along \(x\), the plate collides with it. A
binary variable encodes the opening side, with \(0\) meaning the plate
opens upward (toward row \(r-1\)) and \(1\) downward (toward row \(r+1\)):

\vspace{0.5cm}
\begin{lstlisting}
array[1..n_chunks] of var 0..1: plate;
\end{lstlisting}
\vspace{0.5cm}

An auxiliary variable represents the target row swept by each plate:

\vspace{0.5cm}
\begin{lstlisting}
array[1..n_chunks] of var 0..boxes_along_y+1: target;
constraint forall(g in 1..n_chunks)(
  target[g] = y_row[g] - 1 + 2*plate[g]
);
\end{lstlisting}
\vspace{0.5cm}

When \texttt{plate[g] = 0} the expression yields \texttt{y\_row[g] - 1};
when \texttt{plate[g] = 1} it yields \texttt{y\_row[g] + 1}. The domain
is widened to \([0,\ \texttt{boxes\_along\_y}+1]\) so that the boundary
values \(0\) and \(\texttt{boxes\_along\_y}+1\) represent ``off the
pallet'': they never match a real row index and therefore encode a free
edge without any special case.

The no-collision rule then follows. Because release order coincides with
belt order, \texttt{pred\_chunk} is already on the pallet whenever
\texttt{pred\_chunk} $<$ \texttt{current\_chunk}. The range
\texttt{pred\_chunk in 1..current\_chunk-1} is a compile-time parameter,
so the solver generates only the relevant pairs:

\vspace{0.5cm}
\begin{lstlisting}
constraint forall(curr_chunk in 2..n_chunks,
                  pred_chunk in 1..curr_chunk-1)
(
  not(
       y_row[pred_chunk] = target[curr_chunk]
    /\ x_start[pred_chunk] < x_start[curr_chunk] + chunk_len[curr_chunk]
    /\ x_start[curr_chunk] < x_start[pred_chunk] + chunk_len[pred_chunk]
  )
);
\end{lstlisting}
\vspace{0.5cm}

The body is a direct negation of a single conjunction:
\texttt{pred\_chunk} cannot share the target row of
\texttt{curr\_chunk} while overlapping it along \(x\).